\subsection{The nonflat disk defect survives successive saturated free kernels}
\label{arch136:nonflat-syzygy}

\paragraph{Archive and dependency scope.}
The application is to cote136, physical p.~117/footer116/scan117
and physical p.~121/footer120/scan121. This refines the test of
the bottom parameter module in a strict coherent Moore
realization. It uses ARCH-136-NONFLAT-TAIL-DISK only for its
explicit closed-element calculation and pure divided-power
detector; both arguments are restated where needed. All other
inputs are elementary, tagged none.

Put \(R=\mathbb Z_{(p)}\), \(S=\Gamma_R(T)\), and write
\(S\langle w\rangle\) for a free coefficient module with its
generator in weight \(w\). For an \(S\)-module \(M\), use
the free even disk
\[
       \mathscr D_S(M)=\Gamma_S(M[0])\otimes_S\Lambda_S(M[1]),
                         d x_m=e_m .                    \tag{NS.1}
\]
It retains the coefficient PD structure and uses the special
divided-power differential.

\paragraph{ARCH-136-NONFLAT-SYZYGY-DISK.}
There are weightwise \(R\)-free modules \(P,Q,C\), with
\(P\) free over \(S\), and exact inclusions
\[
            Q\subset P,\qquad P/Q=C\subset F,
                                                               \tag{NS.2}
\]
where \(F\) is free and both displayed inclusions are
saturated with respect to every nonzero coefficient of \(S\),
such that
\[
                         M_0=P+p^{-1}Q\subset P[1/p]
                                                               \tag{NS.3}
\]
has a free even disk with nonzero order-\(p\) cohomology in
form degree zero in an unbounded sequence of weights.
All generators of the ambient free modules can be required
to have weight at least any prescribed positive bound, and
all weight components can be kept finite free over \(R\).
Moreover the construction can be placed in a chain of any
prescribed finite number of successive saturated free kernels.

These assertions refute the indicated structural hypotheses
as sufficient conditions for integral even-disk exactness.
They do not identify the particular higher syzygy of \(I_M\)
occurring in every chosen coefficient-tail resolution.

\paragraph{A two-generator saturated kernel.}
Fix \(h=p^a\), and put
\[
 x=T^{[h]},\qquad y=T^{[ph]},\qquad
 c=\frac{(ph)!}{(h!)^p}=p\kappa,\qquad \kappa\in R^\times.
                                                               \tag{NS.4}
\]
Thus \(x^p=cy\). Consider
\[
 C_h=\ker\bigl(S e_0\oplus S e_1\longrightarrow S\bigr),
 \quad |e_0|=0,\quad |e_1|=h,
 \quad e_0\mapsto p,\quad e_1\mapsto x.
                                                               \tag{NS.5}
\]
The image is the ideal \((p,x)\), a torsion-free submodule
of the domain \(S\). Hence \(C_h\) is \(S\)-saturated in
the displayed free module. It is generated by
\[
        u=(-x,p),\qquad
        v=(-y,\kappa^{-1}x^{p-1}),
        \qquad |u|=h,\quad |v|=ph.                     \tag{NS.6}
\]
To check generation, reduce \(p a+x b=0\) modulo \(p\).
The digit decomposition of \(S/pS\) identifies \(\bar x\)
with one truncated-polynomial generator, so
\(\operatorname{Ann}(\bar x)=(\bar x^{p-1})\).
Thus \(b=p r+x^{p-1}s\). It follows that
\(a=-xr-\kappa y s\), and
\((a,b)=r u+\kappa s v\), proving the claim.

These generators satisfy the exact identities
\[
                     yu=xv,\qquad x^{p-1}u=cv .          \tag{NS.7}
\]
Accordingly the element
\[
                  K_h=y^{p-1}\gamma_p(x_u)
                                  -c\gamma_p(x_v)       \tag{NS.8}
\]
is a closed zero-form in (NS.1) for \(M=C_h\), in weight
\(p^2h\) and arity \(p\). Indeed apply \(\gamma_{p-1}\)
to the first identity in (NS.7), then use the second identity
in the form-one module to obtain
\[
 y^{p-1}\gamma_{p-1}(x_u)e_u
   =x^{p-1}\gamma_{p-1}(x_v)e_u
   =c\gamma_{p-1}(x_v)e_v.
\]
Multiplying that same divided-power identity by \(x_u\)
instead gives \(pK_h=0\).

\paragraph{Nonvanishing for this actual kernel.}
Let \(\mathfrak n=(p,S_+)\). Every relation
\(A u+B v=0\) has \(A,B\in\mathfrak n\).
For the second component of such a relation is
\[
                    pA+\kappa^{-1}x^{p-1}B=0.
\]
Its weight-zero term gives \(A_0=0\). The coefficient of
weight \((p-1)h\) then shows \(B_0\in pR\), since
\(x^{p-1}\) is a \(p\)-unit multiple of
\(T^{[(p-1)h]}\). This proves the assertion for all relations,
including nonhomogeneous ones by taking their weight components.

Use the coefficient ring map
\[
 S\longrightarrow D=\mathbb F_p[X,Y]/(X^p,Y^p)
\]
which retains the digit generators \(T^{[h]},T^{[ph]}\)
as \(X,Y\), and kills all the other digit generators.
It sends \(\mathfrak n\) into \(\mathfrak m=(X,Y)\),
and every element of \(\mathfrak m\) has zero \(p\)-th power.
After this base change, the presentation with generators
\(u,v\) still has all relation coefficients in
\(\mathfrak m\); right exactness of tensoring is sufficient
for this assertion.

In arity \(p\), projection of the free module divided-power
algebra onto its two pure generators \(\gamma_p(u),\gamma_p(v)\)
therefore descends to
\[
                   \Gamma_D^p(C_h\otimes_S D)
                                   \longrightarrow D\oplus D.
                                                               \tag{NS.9}
\]
Here products of positive arities have zero pure projection,
because their pure coefficients contain \(\binom pi=0\)
for \(0<i<p\). The pure projection of the \(p\)-th divided
power of any linear relation is also zero: it consists of
the \(p\)-th powers of its coefficients in \(\mathfrak m\).
All other arity-\(p\) PD-ideal relations are such products,
exactly as in ARCH-136-NONFLAT-TAIL-DISK.
This is a detector using the underlying relative module
divided-power functor and a ring homomorphism; no preservation
of the coefficient PD operations by \(S\to D\) is asserted.

The image of (NS.8) in (NS.9) is
\((Y^{p-1},0)\), which is nonzero. Thus \(K_h\neq0\).
It is already a nonzero cohomology class, since there are no
negative forms, and its exact additive order is \(p\).

\paragraph{The bounded lattice with two successive kernels.}
Take the free cover
\[
 P_h=S f_u\oplus S f_v\xrightarrow{\epsilon}C_h,
 \qquad f_u\mapsto u,\quad f_v\mapsto v,
 \qquad Q_h=\ker\epsilon,
\]
with the indicated generator weights \(h,ph\), and set
\(M_h=P_h+p^{-1}Q_h\subset P_h[1/p]\).
The kernel \(Q_h\) is \(S\)-saturated because its quotient
is \(C_h\subset S e_0\oplus S e_1\).
The quotient \(C_h\) itself is a saturated kernel by (NS.5).
Thus these are precisely two successive saturated free kernels,
not merely an arbitrary torsion-free ideal in a free module.

Inside \(M_h\), put
\[
 \widetilde u=f_u,
 \qquad
 \widetilde v=f_v+
             \frac{x^{p-1}f_u-c f_v}{c}
             =\frac{x^{p-1}}c f_u .                     \tag{NS.10}
\]
The numerator in the fraction belongs to \(Q_h\), and
\(v_p(c)=1\), so the expression belongs to \(M_h\).
The relations (NS.7) hold for the tilded elements as well.
The resulting closed order-dividing-\(p\) expression maps
under \(\epsilon:M_h\to C_h\) to the nonzero (NS.8),
and so has exact order \(p\). This proves the required
failure for a bounded lattice of the stated form.

To obtain a single example with unbounded weights, take the
direct sum of these constructions over \(h=p^a\) above a
prescribed bound, shifting the weight of every module in
the \(h\)-th block by \(h\). The free ambient generators
then all have weight at least \(h\); only finitely many
blocks contribute to any fixed weight. Direct sums of the
displayed kernel sequences are exact, with torsion-free
quotients, so both saturation conditions persist.
Every weight component of the bounded lattice is a submodule
of a finite free \(R\)-module and is therefore free over \(R\).
The classes (NS.8) in the shifted blocks have weights
\(p^2h+ph\), and remain nonzero in the direct sum module's
PD algebra by the projection back to their block. Their
weights are unbounded. This proves all uniform assertions
in the statement.

\paragraph{Arbitrarily long saturated kernel chains.}
The same calculation gives an exact elementary syzygy formula.
For \(1\leq j<p\), put \(J_j=(p,x^j)\subset S\).
Then
\[
 0\longrightarrow J_{p-j}\langle jh\rangle
 \longrightarrow S\oplus S\langle jh\rangle
 \xrightarrow{(p,x^j)}J_j\longrightarrow0.              \tag{NS.11}
\]
In fact \(\operatorname{Ann}_{S/pS}(\bar x^j)
=(\bar x^{p-j})\). Writing the second coordinate as
\(p r+x^{p-j}s\) gives the kernel generators
\((-x^j,p)\) and \((-\kappa y,x^{p-j})\);
projection to that coordinate identifies the kernel with
\(J_{p-j}\), with the displayed weight shift.
All the ideals \(J_j\) are torsion-free over the domain
\(S\). Thus every kernel in (NS.11) is saturated.
Iterating alternates \(J_j\) and \(J_{p-j}\), and supplies
any prescribed finite number of successive saturated free
kernel presentations. Since \(C_h\simeq J_{p-1}\langle h\rangle\),
it can be placed in such a chain. A sufficiently large common
weight shift makes every free generator in any chosen finite
chain positive and above a prescribed bound. Performing these
shifts blockwise as above preserves unbounded weights of the
detected classes.

\paragraph{Exact limitation for the active construction.}
Neither coefficient-torsion-freeness, integer saturation,
two successive saturated free kernels, nor any prescribed
finite kernel depth yields a general even-disk Poincare lemma.
The class just exhibited has an explicit detector, so this is
not an inference from nonflatness alone.
The argument does not prove that the particular syzygy of
the coefficient tail \(I_M\) in every chosen projective
resolution contains these modules or these detected relations.
That more specific test, or a subsequent low-degree quotient
and repair with its own admissibility proof, remains separate.

\paragraph{Verification status.}
The author has read the complete saved proof, including kernel
generation, minimal relation coefficients, the pure-PD detector,
both saturated kernels, bounded lattices, direct-sum weight bounds,
and the alternating syzygy formula. The parent has independently read
the complete saved proof and communicated PASS on 2026-11-22.
Reviewer1348 has independently read the complete saved proof and
communicated PASS as well. The separate actual-tail-resolution
theorem is not part of these checks.
